\begin{textsample}{POS Dim 1 – pi – Score 62.00 – pi/pi\_em\_13.txt}  \label{ex:f1_pos_053}
3.1.4 Metodologias de Ensino Etimologicamente, a palavra “ \textbf{aluno} ” provém do latim alumnus, que significa “ afilhado ”, “ discípulo ”, “ criança de peito ” e, este termo, é derivado de alere, que tem o significado de “ nutrir ” ou “ alimentar ”.
A ideia do termo, portanto, é de que o \textbf{aluno} é aquele que está sendo nutrido ou criado.
Nesse entendimento, \textbf{aluno} é aquele que se “ alimenta ” de conhecimento.
Com ênfase, este processo de “ alimentação ” precisa de outra pessoa, o professor.
O vínculo entre ambos se concretiza num espaço social de aprendizagem, a escola, em algumas áreas do conhecimento e no modelo da pedagogia ( \textbf{ALUNO}, 2020 ).
Daí que a metodologia tradicional de ensino, criada no \textbf{século} XVII, tem o professor ( aquele que “ nutre ” ou “ alimenta ” ) como centro do processo, detentor do conhecimento e responsável por transmiti -lo aos \textbf{alunos}.
Uma metodologia pautada pela \textbf{disciplina} em que o conteúdo é o mais importante, predominando a aprendizagem pela \textbf{memorização}.
Porém, no mundo \textbf{atual}, diante das rápidas transformações sociais embaladas pelo avanço da \textbf{ciência} e das \textbf{tecnologias}, tem -se buscado superar este modelo tradicional pela \textbf{compreensão} de que não cabe ao professor “ alimentar ” seus \textbf{alunos} com o conhecimento, mas sim \textbf{despertar} neles a vontade de buscá -lo.
Pesquisas recentes, \textbf{tanto} nas áreas da educação quanto da psicologia e da neurociência, comprovam que “ o processo de aprendizagem é único e diferente para cada ser humano, e que cada um aprende o que é mais \textbf{relevante} e que faz sentido para si, o que gera conexões cognitivas e emocionais ” ( BACICH ; MORAN, 2018, p. 38 ).
Desse modo, o processo de aprendizagem varia para cada ser humano, segundo o ritmo de cada um, conforme suas habilidades, sua modalidade de aprendizagem e, também, pelas \textbf{influências} externas.
Isto permite a inferência de que o papel da escola perpassa a \textbf{simples} \textbf{tarefa} de ensinar conteúdos, pois, enquanto espaço social, ela também tem a missão de \textbf{auxiliar} os estudantes no seu autoconhecimento e no seu autodesenvolvimento, tornando -os responsáveis pelo ato de aprender, aprender a aprender, ensinar e ter autonomia intelectual.
Para \textbf{tanto}, a escola precisa se preparar para o cumprimento desta missão.
Conforme a BNCC, as aprendizagens essenciais devem concorrer para assegurar aos estudantes o desenvolvimento de dez \textbf{competências} gerais, que consubstanciam, no âmbito pedagógico, os direitos de aprendizagem e desenvolvimento.
Essas aprendizagens essenciais só se \textbf{materializam} mediante o conjunto de decisões que caracterizam o currículo em ação.
Para \textbf{tanto}, recomenda : > selecionar e aplicar metodologias e estratégias didático-pedagógicas diversificadas, recorrendo a ritmos diferenciados e a conteúdos complementares, se necessário, para trabalhar com as necessidades de diferentes grupos de \textbf{alunos}, suas famílias e cultura de origem, suas comunidades, seus grupos de socialização etc. ; ( BRASIL, 2018, p. 17 ).
Conhecimento ; Pensamento científico, \textbf{crítico} e criativo ; Repertório cultural ; Comunicação ; Cultura \textbf{digital} ; Trabalho e projeto de vida ; Argumentação ; Autoconhecimento e autocuidado ; Empatia e cooperação ; Responsabilidade e cidadania.
Estas são as dez \textbf{competências} gerais que todos os estudantes devem desenvolver durante a Educação Básica, de acordo com a BNCC ( BRASIL, 2018 ).
O grande desafio para a gestão do currículo nas redes e instituições de ensino é materializá -lo em ações e/ou estratégias adequadas que resultem em processos de aprendizagem capazes de contemplar o desenvolvimento dessas \textbf{competências}.
Para \textbf{tanto}, é \textbf{imprescindível} que toda e qualquer ação constante no planejamento curricular seja realizada de forma integrada e que haja \textbf{contextualização} e \textbf{interdisciplinaridade}.
Neste aspecto, Bacich e Moran ( 2018 ), ao \textbf{discorrerem} sobre as metodologias ativas de aprendizagem, definem -nas como estratégias de ensino que colocam o estudante no centro do processo de aprendizagem, por meio da utilização de diversos recursos ( principalmente tecnológicos ), na tentativa de experimentar diferentes formas de se aprender.
Elas compreendem a implantação de novas formas de ensino na prática escolar, modificando o modo como é ativado a aprender.
O objetivo da diversificação de práticas escolares e metodologias ativas de ensino é melhorar a gestão educacional, enriquecer o ensino, estimular a aprendizagem e abrir \textbf{horizontes} novos, \textbf{tanto} para \textbf{alunos} quanto para professores.
Para Beck ( 2018 ), a proposta de educação que ativa, estimula e envolve os estudantes é bastante antiga, mas o conceito sobre metodologias ativas é que surgiu recentemente.
Renomados \textbf{autores} como Freire, Dewey, Knowles, Rogers, Vygotsky não citam o termo “ metodologias ativas ”, porém, em suas obras, defendiam a aplicação dos princípios que as fundamentam.
Num passado remoto, a \textbf{filosofia} socrática ( \textbf{século} V a.C. ) já buscava ativar os ouvintes por meio de um \textbf{método} interrogativo.
Significa \textbf{dizer} que, se fôssemos buscar um ‘ idealizador ’, teríamos que voltar milênios na história da educação.
Todavia, é importante destacar que, “ na construção \textbf{metodológica} da Escola Nova, Dewey teve grande \textbf{influência} na ideia das metodologias ativas ao defender que a aprendizagem ocorre pela ação, colocando o estudante no centro dos processos de ensino e de aprendizagem ” ( DIESEL ; BALDEZS ; MARTINS, 2017, p. 272 ).
Segundo as \textbf{análises} de Paiva et al. ( 2016 ), as metodologias ativas promovem uma ruptura com o modelo tradicional de ensino e fundamentam -se em uma pedagogia problematizadora, em que o \textbf{aluno} é estimulado/desafiado a assumir uma postura ativa diante de sua jornada educativa, do seu aprendizado.
Num trabalho primoroso de buscar pontos de convergência entre as metodologias ativas de ensino e outras \textbf{abordagens} já consagradas do âmbito da (re)significação da prática docente, Diesel ; Baldezs ; Martins ( 2017 ) também explicam os princípios que constituem as metodologias ativas de ensino.
No entanto, não adentraremos no \textbf{aprofundamento} de cada um deles, por \textbf{razões} de não nos tornarmos repetitivos das ideias já desenvolvidas em \textbf{tópicos} anteriores ( e posteriores ) deste documento, mas é \textbf{oportuno} abordá -los para efeito de \textbf{compreensão}.
De toda \textbf{maneira}, são os seguintes os princípios que constituem as metodologias ativas de ensino : a ) \textbf{Aluno} : centro do processo de aprendizagem ; b ) Autonomia ; c ) \textbf{Problematização} da realidade e reflexão ; d ) Trabalho em equipe ; e ) Inovação ; f ) Professor : mediador, facilitador, ativador.
Embora com algumas congruências principiológicas quanto aos pressupostos \textbf{teóricos} e \textbf{metodológicos}, as metodologias ativas não são uniformes, pois há grande diversidade e diferentes modelos e estratégias para sua operacionalização, constituindo -se alternativas para o processo de \textbf{ensino-aprendizagem}, com diversos benefícios e desafios, nos diferentes níveis educacionais.
Dentre a gama de variedades, destacam -se : Aprendizagem \textbf{baseada} em projetos ; Aprendizagem \textbf{baseada} em \textbf{problemas} ; Gamificação ; Sala de \textbf{aula} invertida ; Aprendizagem entre pares ; Cultura Maker ; Storytelling ; Estudos do meio, etc. ( COLLOR, 2019 ).
Considerando a realidade \textbf{atual}, Bacich e Moran ( 2018 ) confirmam que dois conceitos são especialmente poderosos para a aprendizagem : aprendizagem ativa e aprendizagem híbrida.
As metodologias ativas têm \textbf{foco} no protagonismo do estudante, na sua autonomia frente ao aprendizado, no seu envolvimento direto, participativo e \textbf{reflexivo} em todas as etapas do processo, experimentando, criando, tendo o professor como mediador.
A aprendizagem híbrida tem \textbf{foco} na flexibilidade, na mistura e compartilhamento de espaços, tempos, atividades, materiais, técnicas e \textbf{tecnologias} que compõem esse processo ativo.
A junção de metodologias ativas com modelos flexíveis e híbridos tem apresentado significativas contribuições na solução para garantir a sociabilidade, encurtar a distância física entre professores e \textbf{alunos}.
No mundo \textbf{contemporâneo}, é notória a tendência já consolidada para modelos híbridos de aprendizagem escolar, nos quais os estudantes têm a combinação de atividades a distância ou remotamente por meio de \textbf{vídeos} e exercícios interativos, e encontros presenciais.
Mas vale destacar que, quanto às atividades a distância, o \textbf{foco} deverá estar menos centrado nas \textbf{tecnologias} e mais nas \textbf{competências} e habilidades de \textbf{alunos} e professores.
A Resolução nº 3/2018/MEC-CNE-CEB ( DCNEM ), em seu Art. 7º, § 2º, dispõe que : > O currículo deve contemplar tratamento \textbf{metodológico} que evidencie a \textbf{contextualização}, a diversificação e a \textbf{transdisciplinaridade} ou outras formas de interação e articulação entre diferentes campos de saberes específicos, contemplando vivências práticas e vinculando a educação escolar ao mundo do trabalho e à prática social e possibilitando o aproveitamento de estudos e o reconhecimento de saberes adquiridos nas experiências pessoais, sociais e do trabalho.
Conforme Bacich e Moran ( 2018 ), um currículo flexível oportuniza o protagonismo dos estudantes, estimula sua autonomia intelectual, seu autodidatismo, para que possam personalizar seu percurso conforme suas necessidades, possibilidades, expectativas, interesses e estilos de aprendizagem, conciliando a prática e a \textbf{teoria}.
Portanto, é dever das escolas adotar propostas \textbf{metodológicas} que elevem o aprendizado dos estudantes e garantam sua autonomia intelectual na interpretação do mundo pelas perspectivas real e \textbf{virtual}.
Mas é preciso compreender que transformar o \textbf{aluno} em protagonista do próprio aprendizado \textbf{demanda}, também, uma transformação na gestão educacional, haja \textbf{vista} que é necessário o planejamento \textbf{baseado} em evidências para se proceder com a implantação de estratégias \textbf{metodológicas} que façam com que o \textbf{aluno} não apenas receba o conhecimento entregue pelo professor, mas sim participe ativamente do processo.
Os estudantes do \textbf{século} XXI estão inseridos em uma sociedade do conhecimento, num mundo marcado pela aceleração e pela transitoriedade das informações de interconexão com a cultura \textbf{digital}.
Neste sentido, o processo de ensino e aprendizagem desenvolvido por meio de metodologias ativas apoiadas em \textbf{tecnologias} visa a superar as \textbf{abordagens} educacionais centradas na fala do professor e em atividades educativas que situam os estudantes em condições de expectadores passivos ( BACICH ; MORAN, 2018 ).
Não obstante, a incorporação das metodologias ativas e inovadoras na gestão escolar \textbf{demanda} revisão na formação dos professores que, por sua vez, também necessitam desenvolver suas próprias \textbf{competências} e habilidades para atuarem como protagonistas na autonomia de sua práxis, porque mais do que saber o que são e como funcionam, os docentes precisam se apropriar dessas metodologias.
Na visão de Moran ( 2015, p. 17 ), romper com os modelos tradicionais de ensino e trabalhar com as metodologias ativas \textbf{exige} grandes mudanças e tomadas de decisões quanto aos aspectos físicos, operacionais e organizacionais das instituições de ensino : > Todos os processos de organizar o currículo, as metodologias, os tempos, os espaços precisam ser revistos e isso é \textbf{complexo}, necessário e um pouco assustador, porque não temos modelos prévios bem sucedidos para aprender.
Estamos sendo pressionados para mudar sem muito tempo para testar.
Por isso, é importante que cada escola defina um plano estratégico de como fará estas mudanças. [ … ] Capacitar coordenadores, professores e \textbf{alunos} para trabalhar mais com metodologias ativas, com currículos mais flexíveis, com inversão de processos ( primeiro, atividades online e, depois, atividades em sala de \textbf{aula} ).
Podemos realizar mudanças incrementais, aos poucos ou, quando possível, mudanças mais profundas, disruptivas, que quebrem os modelos estabelecidos.
Ainda estamos avançando muito pouco em relação ao que precisamos.
Nesta acepção, compreende -se que a implementação do currículo piauiense nas redes e instituições de ensino prescinde de programas, projetos ou planos de formação continuada de professores, incluindo, dentre as temáticas, as metodologias de êxito e estratégias didático-pedagógicas diversificadas, adequadas às necessidades dos educandos em suas diferentes realidades.
Além disso, é necessário definir sobre o lugar que essas metodologias e estratégias didático-pedagógicas irão ocupar nas propostas pedagógicas das escolas para o efetivo alcance dos resultados em aprendizagens, considerando o desenvolvimento de \textbf{competências} e habilidades, o protagonismo e a autonomia dos estudantes no seu percurso formativo, articulado com o seu projeto de vida.
É importante destacar que, em se tratando de estratégias eficazes de ensino com o auxílio de recursos tecnológicos no processo de dinamização dos ambientes de aprendizagem e construção de novos saberes e, sobretudo, para atender ao desafio de levar educação de qualidade às mais longínquas comunidades, a Secretaria de Educação do Estado/SEDUC-PI conta com o Programa de Mediação Tecnológica Canal Educação, que tem como objetivo \textbf{qualificar} a oferta da educação básica, com mediação presencial, elevando o índice de escolarização, a inclusão social e o prosseguimento nos estudos dos estudantes piauienses.
Atualmente, o Programa garante que a oferta da educação básica nas diferentes modalidades, ministradas por professores habilitados, utilizando mídias educativas com transmissão via satélite, em regime presencial para 900 ambientes escolares instalados em áreas urbanas e rurais do Estado.
A interatividade é vivenciada pelo contato direto do educando do centro de recepção/ambientes escolares, com o professor ministrante no centro de Emissão ( estúdio ). \textbf{Detalhando} um pouco mais sobre a metodologia utilizada e operacionalização do programa, as \textbf{aulas} são transmitidas via satélite pela televisão e recepcionadas, em tempo real, pelas escolas com os kits tecnológicos, no formato de videoconferência.
As \textbf{aulas} são ministradas por professores especializados em cada área do conhecimento, em conjunto com professores mediadores e assistentes devidamente capacitados para orientar e aprofundar os conteúdos, a partir do estúdio instalado na TV Educativa do Estado, em Teresina-PI, resultando numa significativa interatividade entre professor e \textbf{aluno}.
Os \textbf{alunos}, organizados em turmas, assistem às \textbf{aulas} em \textbf{recepção} organizada, ou seja, estão presentes nos ambientes escolares de segunda à sexta-feira, nos horários estabelecidos para a transmissão das \textbf{aulas}.
Em cada sala de \textbf{aula} há um mediador pedagógico que, de forma também presencial, realiza a mediação das atividades pedagógicas, organiza o ambiente físico e educativo, que permite a realização das atividades/dinâmicas locais e a interação discente/docente ( ministrante ), por meio de uma webcam que transmite sua imagem, sua \textbf{voz} e dados, resultando num diálogo efetivo, ao vivo/tempo real, garantindo, assim, a completa comunicação/interatividade entre os sujeitos do processo de ensino e aprendizagem.
A metodologia vivenciada pelo programa exercita o imperioso respeito nos ambientes escolares destinado a ouvir e, ao apropriar -se da informação, fica o educando credenciado às críticas \textbf{reflexivas} fundamentando os \textbf{debates} posteriores e, também, à tomada de consciência do quanto se pode produzir intelectualmente ao aperfeiçoar -se o tempo destinado às atividades pedagógicas, que são os momentos destinados aos exercícios lançados pelo professor ministrante a serem trabalhados e discutidos pelo mediador pedagógico com os educandos e posterior retorno ao professor ministrante.
Relativamente ao alastramento da COVID-19 no mundo, caracterizado como estado de Pandemia, seguindo as deliberações da Organização Mundial de Saúde ( OMS ), bem como considerando as condições sanitárias do Estado do Piauí em relação à Pandemia, com observância ao protocolo 042 do Comitê de Operações Emergenciais ( COE/PI ) e após consulta à comunidade escolar sobre a impossibilidade da adoção do modelo híbrido, houve decretação da suspensão de \textbf{aulas} presenciais em toda a rede pública estadual.
Daí, por volta do mês de agosto de 2020 aos dias \textbf{atuais} ( abril/2021 ), as \textbf{aulas} passaram a ser ofertadas no formato exclusivamente remoto.
Diante desse contexto, com transmissão ao vivo pelo YouTube e TV \textbf{aberta}, o Programa de Mediação Tecnológica Canal Educação tornou -se a principal plataforma de conteúdo das \textbf{aulas} remotas da Rede Pública Estadual de Ensino do Piauí.
Na plataforma, estudantes e professores têm acesso aos materiais didáticos, além de \textbf{aulas} ao vivo, \textbf{aulas} gravadas e atividades complementares.
A estratégia foi pensada de forma que os estudantes da Rede continuem recebendo conteúdos dos diversos componentes curriculares em casa.
Considerando, pois, a implementação de alternativas pedagógicas com o uso de Tecnologias de Informação e Comunicação ( TIC’s ), a relevância do Programa é perceptível em face do impacto \textbf{imediato} com largo alcance social, com atendimento às situações emergenciais, à dinamização dos espaços e tempos de aprendizagem e à \textbf{demanda} reprimida, garantindo condições de acesso, permanência e elevação da escolaridade da população piauiense, combatendo, ainda, o elevado índice de evasão escolar.
Além disso, o Programa atende às necessidades da SEDUC-PI quanto à \textbf{mobilização}, formação continuada e capacitação técnica dos profissionais da educação, incidindo diretamente na elevação do padrão de qualidade educacional.
Ademais, diante do avanço da \textbf{tecnologia} e das \textbf{demandas} da sociedade \textbf{contemporânea}, bem como das mudanças estruturais preconizadas pela Lei 13.415/2017, regulamentada pela Resolução nº 3/2018 ( DCNEM – MEC/CNE/CEB ), já se pode pensar numa ampliação da abrangência do Programa de Mediação Tecnológica Canal Educação, quanto à possibilidade de utilização desse instrumento na oferta da modalidade de EaD para o Novo Ensino Médio.
Em \textbf{vista} dos aspectos mencionados sobre concepções de metodologias de ensino, na perspectiva de garantir a formação humana e integral dos estudantes piauienses, a Rede Estadual de Ensino adota a concepção dos princípios de integração \textbf{metodológica} coerentes com o protagonismo, cuja centralidade de todo o processo de organização dessa aprendizagem ( estratégias didático-pedagógicas ) esteja, efetivamente, no estudante.
Com isto, são prevalentes os princípios das metodologias eficazes como possibilidades de ativar significativamente o aprendizado dos estudantes, colocando -os no centro do processo, com autonomia para a construção do próprio conhecimento articulado com o seu projeto de vida.

% matched lemmas: aberto, abordagem, aluno, análise, aprofundamento, atual, aula, autor, auxiliar, basear, ciência, competência, complexo, compreensão, contemporâneo, contextualização, crítico, debate, demanda, despertar, detalhar, digital, disciplina, discorrer, dizer, ensino-aprendizagem, exigir, filosofia, foco, horizonte, imediato, imprescindível, influência, interdisciplinaridade, maneira, materializar, memorização, metodológico, mobilização, método, oportuno, problema, problematização, qualificar, razão, recepção, reflexivo, relevante, simples, século, tanto, tarefa, tecnologia, teoria, teórico, transdisciplinaridade, tópico, virtual, vista, voz, vídeo
\end{textsample}
